\section{Scoring candidates under restrictions}
\label{sec:selection}

Applied bounds are computed under substantive restrictions, so the selection rule must handle a polyhedral $\restrset$, where widths
are finite and the dichotomy of \cref{prop:dichotomy} no longer holds. The
key structural fact is that the candidate's entire effect on the
feasible set passes through two scalars.

\begin{proposition}[Adding a site adds two equality constraints]
\label{prop:refinement}
Suppose \cref{ass:model} holds for the pool and for candidate $k$, and
the eigenvalue condition holds for the augmented pool. Let
$\feasset'$ denote the feasible set \eqref{eq:feasible-set} built from
the augmented moments. If $\cand{k}\notin\affhull(\pool)$, then, with
$\newdir{k}$ normalized to unit length,
\begin{equation}
  \feasset'
  =\feasset\cap
  \bigl\{\beta:\ \newdir{k}^\top\beta_a=\reveal{a}{k},\ a\in\{0,1\}\bigr\},
  \qquad
  \reveal{a}{k}=\newdir{k}^\top\beta_a,
  \label{eq:refinement}
\end{equation}
where $\reveal{a}{k}$ is identified by the augmented pool but not by
$\pool$. If $\cand{k}\in\affhull(\pool)$, then $\feasset'=\feasset$.
The bounds \eqref{eq:endpoint-programs} built from the augmented moments
equal the bounds built from the base objective $(\bT,\ellvec)$ over
$\feasset'$: the changes in $\bT$ and $\ellvec$ induced by re-pooling
shift both endpoints by the same amount tied to the re-decomposition of
$\psi_T$, and the target effect interval is unchanged by them.
\end{proposition}

\cref{prop:refinement} is the basis of the design rule. The direction
$\newdir{k}$ is computable from the profiles before the experiment. The scalars
$\reveal{0}{k}$ and $\reveal{1}{k}$ are what the candidate experiment
reveals: the projection of each arm's fixed effects on the new
direction. Everything else about the candidate's data (its sample size, its noise
level, and its covariate mix) affects only how precisely those two
scalars are estimated. It does not affect the identified set.

\paragraph{The predicted interval.}
Define the design-time width map
\begin{equation}
  \width_k(\ct;\,y_0,y_1)
  =\text{width of \eqref{eq:endpoint-programs} over }
  \feasset\cap\{\newdir{k}^\top\beta_a=y_a,\ a\in\{0,1\}\},
  \label{eq:width-map}
\end{equation}
a pair of linear programs the investigator can solve before the
experiment, for any conjectured $(y_0,y_1)$. Before the experiment we evaluate it at a
forecast. The neutral forecast extends the pool's own fit. The minimum-norm
solution $\what\beta_a$ of the estimated equality system sets the
unidentified coordinates to zero, so its forecast is
$\what y_a=\newdir{k}^\top\what\beta_a=0$. Any substantive prior that
maps the candidate's context to expected outcomes can replace it. We
call $\width_k(\ct;\what y_0,\what y_1)$ the predicted width and use it,
maximized over $\targset$, as the score in \eqref{eq:minimax}.

\paragraph{Pricing the forecast error.}
The prediction is wrong only to the extent that the forecast is wrong.
The sensitivity follows from a standard envelope argument. At a bound
with a unique optimizer and dual, the derivative of the bound in $y_a$
is the multiplier on the added row, so for small errors
\begin{equation}
  v_j(y)\approx v_j(\what y)
  +\sum_{a\in\{0,1\}}\lambda^{u}_{j,a}\,(y_a-\what y_a),
  \qquad j\in\{-,+\},
  \label{eq:dual-pricing}
\end{equation}
where $\lambda^{u}_{j,a}$ is bound $j$'s dual variable on the constraint
$\newdir{k}^\top\beta_a=y_a$. This is the same envelope argument that the inference theory of the
companion paper uses, applied to the added rows. An investigator who can bound the forecast error, say
$|y_a-\what y_a|\leq B$, obtains a range of predicted intervals by
adding $\sum_a(|\lambda^{u}_{-,a}|+|\lambda^{u}_{+,a}|)\,B$ to the
predicted width, or exactly by re-solving \eqref{eq:width-map} on a grid
of $y$ values, since the map is piecewise linear in $y$. In our
simulations and in the case study the ranking of candidates is stable
across these variants, because the ranking depends on which direction is resolved and not on
the forecast level.

\paragraph{Three kinds of candidate.}
The three categories from the companion paper's leave-one-out exercise
also describe what a candidate can do.
\begin{itemize}
\item \emph{Identification.} If $\cand{k}\notin\affhull(\pool)$, the
  candidate adds the equality constraints \eqref{eq:refinement}. This is
  the only way a candidate changes the identified set, and it is
  checkable from profiles alone.
\item \emph{Precision.} If $\cand{k}\in\affhull(\pool)$, the identified
  set is unchanged, and the candidate's value is statistical. More context
variation inside $\Vspan$ tightens the estimates and the standard errors
of the same bounds. Interior sites add precision but not identification.
\item \emph{Feasibility.} A candidate can also restore or destroy
  feasibility. If the pool's program is infeasible under $\restrset$,
  no addition restores it, since $\feasset'\subseteq\feasset$ under the
  model. But an addition whose data conflict with the model can make a feasible
program infeasible. An infeasible augmented program then signals that
conflict, not a property of the candidate.
  \cref{sec:monotonicity} takes this up.
\end{itemize}

\paragraph{Scoring procedure.}
The rule we use in \cref{sec:simulation,sec:case-study} has four steps.
\begin{enumerate}
\item For each candidate, test $\cand{k}$ against $\affhull(\pool)$.
\item For each candidate that expands the hull, solve \eqref{eq:width-map}
  at the forecast over $\targset$.
\item Report the minimax score, the resolved direction, and the dual
  prices of the forecast.
\item Break ties, and rank interior candidates, by predicted precision
  at fixed identification.
\end{enumerate}
The whole computation reuses the base-pool estimate $\what\Gam$ and adds two rows per
candidate, so scoring a pool of $K$ candidates costs $K$ pairs of linear
programs per target beyond the base fit.
